\documentclass{article}

\usepackage{amsmath}
\usepackage{amssymb}
\usepackage{placeins}
\usepackage{graphicx}
\usepackage[colorlinks]{hyperref}
\usepackage{longtable}
\usepackage{booktabs}
\usepackage{array}
\usepackage{ragged2e}
\newcolumntype{L}[1]{>{\RaggedRight\arraybackslash}p{#1}}

\title{Production Scheduling -- Short Term -- Intra-Shift Job Sequencing}
\author{LTPlabs -- APP\&HIPO}
\date{}

\begin{document}

\maketitle

\section{Mixed Integer Linear Formulation}

The daily model fixes which operations are processed on each machine and shift,
as well as their production quantities. The present model only decides the
order of those operations within the shift. Consequently, processing times are
fixed and only sequence-dependent setup times vary with the decision.

The shift length is not a feasibility constraint. Any completion after the end
of the shift is reported as an overrun and can be used as feedback to the daily
model.

The model is solved independently for each shift and each connected component
of machines. In the remainder, let shift $s$ and component $\kappa$ be fixed.

\subsection{Job-level notation}

The daily model identifies an operation through the triple $(p,r,o)$. To keep
the sequencing formulation compact, each active operation is represented by a
single job index
\begin{align}
j \equiv (p,r,o).
\end{align}
All product, routing, and operation information remains attached to $j$; the
short index only avoids repeating the full triple in every expression.

\begingroup
\small
\setlength{\LTleft}{0pt}
\setlength{\LTright}{0pt}
\setlength{\tabcolsep}{4pt}
\renewcommand{\arraystretch}{1.15}

\begin{longtable}{@{}L{0.30\textwidth}L{0.64\textwidth}@{}}
\caption{Sets and indices}
\label{tab:seq-sets} \\

\toprule
\textbf{Set or index} & \textbf{Description} \\
\midrule
\endfirsthead

\caption[]{Sets and indices (continued)} \\
\toprule
\textbf{Set or index} & \textbf{Description} \\
\midrule
\endhead

\bottomrule
\endlastfoot

$s$ & Shift being scheduled. Fixed throughout the model; shifts are solved one at a time. \\

$\kappa$ & Connected component of machines within shift $s$, as defined below. Fixed throughout the model; each component is solved independently. \\

$\mathcal{M}_{\kappa}$ & Machines belonging to component $\kappa$ in shift $s$. \\

$\mathcal{J}_{m}$ & Jobs assigned by the daily plan to machine $m \in \mathcal{M}_{\kappa}$ in shift $s$. \\

$\mathcal{J}_{\kappa}=\bigcup_{m\in\mathcal{M}_{\kappa}}\mathcal{J}_{m}$ & Jobs belonging to component $\kappa$. \\

$m(j)$ & Machine to which job $j$ is assigned. \\

$0_m$ & Dummy origin of machine $m$. It represents the setup state carried from the most recent previous shift in which $m$ was active. \\

$\mathcal{N}_{m}=\{0_m\}\cup\mathcal{J}_{m}$ & Nodes available in the sequence of machine $m$. \\

$\mathcal{A}_{m}=\{(i,j): i\in\mathcal{N}_{m},\ j\in\mathcal{J}_{m},\ i\neq j\}$ & Candidate immediate-successor arcs. Arcs enter real jobs, but never return to the dummy origin. \\

$\mathcal{E}$ & Binding precedence edges $(i,j)$ inherited from the daily plan. Edge $(i,j)$ means that job $j$ depends on output produced by job $i$ in the same shift. \\

\end{longtable}
\endgroup

Each component is formed by connecting machines whenever an edge in
$\mathcal{E}$ runs between them. Therefore, no sequencing decision in one
component constrains a decision in another component.

\subsection{Parameters}

The hat accent continues to denote values fixed by the daily-model solution.

\begingroup
\small
\setlength{\LTleft}{0pt}
\setlength{\LTright}{0pt}
\setlength{\tabcolsep}{4pt}
\renewcommand{\arraystretch}{1.15}

\begin{longtable}{@{}L{0.30\textwidth}L{0.64\textwidth}@{}}
\caption{Parameters}
\label{tab:seq-parameters} \\

\toprule
\textbf{Parameter} & \textbf{Description} \\
\midrule
\endfirsthead

\caption[]{Parameters (continued)} \\
\toprule
\textbf{Parameter} & \textbf{Description} \\
\midrule
\endhead

\bottomrule
\endlastfoot

$p_j$ & Fixed processing time of job $j$. \\

$d_{ij}$ & Setup time required when job $j$ immediately follows node $i$ on their shared machine. For $i=0_m$, it is the setup from the state carried into the shift. \\

$\gamma_{m,t}$ & Setup time incurred on machine $m$ to satisfy one unit of setup type $t$. \\

$\hat c_j\in\{0,1\}$ & Equals 1 when job $j$ is a continuation from the previous shift and must be first on its machine. \\

$\hat e_j\in\{0,1\}$ & Equals 1 when job $j$ is the setup state selected by the daily plan for the end of the shift and must be last on its machine, unless $\hat c_j=1$ too, or $j$ has a same-machine binding successor (see A5). \\

$M$ & Valid upper bound used to deactivate a conditional timing constraint. \\

\end{longtable}
\endgroup

Because quantities are fixed, processing time is data rather than a decision:
\begin{align}
p_j
=
\hat{q}_{p,r,o,s}\,\text{standardTime}_{p,r,o},
\qquad j\equiv(p,r,o)\in\mathcal{J}_{\kappa}.
\label{eq:processing-time}
\end{align}

Let $\Theta_{j,m}$ be the setup characteristics required by job $j$ on machine
$m$, and let $t(\theta)$ be the setup type of characteristic
$\theta\in\Theta_{j,m}$. Transitioning from $i$ to $j$ pays for every
characteristic required by $j$ and absent from $i$:
\begin{align}
d_{ij}
=
\sum_{\theta\in\Theta_{j,m}\setminus\Theta_{i,m}}
\gamma_{m,t(\theta)},
\qquad (i,j)\in\mathcal{A}_{m},\quad i\neq 0_m.
\label{eq:setup-distance}
\end{align}
Dropping a characteristic is free, so $d_{ij}$ is directed. For the dummy
origin, $d_{0_mj}$ is evaluated from the state carried into the shift; it is
zero when no prior state exists.

\subsection{Decision variables}

\begingroup
\small
\setlength{\LTleft}{0pt}
\setlength{\LTright}{0pt}
\setlength{\tabcolsep}{4pt}
\renewcommand{\arraystretch}{1.15}

\begin{longtable}{@{}L{0.30\textwidth}L{0.64\textwidth}@{}}
\caption{Decision variables}
\label{tab:seq-decision-variables} \\

\toprule
\textbf{Variable} & \textbf{Definition} \\
\midrule
\endfirsthead

\caption[]{Decision variables (continued)} \\
\toprule
\textbf{Variable} & \textbf{Definition} \\
\midrule
\endhead

\bottomrule
\endlastfoot

$y_{ij}\in\{0,1\}$ & Equals 1 if job $j$ immediately follows node $i$ on their shared machine; defined for $(i,j)\in\mathcal{A}_{m}$. In particular, $y_{0_mj}=1$ means that $j$ is first. \\

$B_j\geq0$ & Start time of job $j$, measured from the beginning of shift $s$. The setup of $j$ begins at this instant. \\

$C_j\geq0$ & Completion time of job $j$, after its setup and processing. \\

$C^{\max}_{\kappa}\geq0$ & Completion time of the last job in component $\kappa$. \\

\end{longtable}
\endgroup

\FloatBarrier

\subsection{Formulation}

\subsubsection{Objective function}

\begin{align}
\textbf{minimize}\qquad C^{\max}_{\kappa}.
\label{eq:objective}
\end{align}

\subsubsection{A. Sequence structure}

\noindent\textbf{A1. One sequence leaves each machine origin}
\begin{align}
\sum_{j\in\mathcal{J}_{m}}y_{0_mj}=1,
\qquad
\forall m\in\mathcal{M}_{\kappa}:\mathcal{J}_{m}\neq\emptyset.
\label{eq:A1-origin}
\end{align}

\noindent\textbf{A2. Every job has exactly one predecessor}
\begin{align}
\sum_{i\in\mathcal{N}_{m}\setminus\{j\}}y_{ij}=1,
\qquad
\forall m\in\mathcal{M}_{\kappa},\quad
\forall j\in\mathcal{J}_{m}.
\label{eq:A2-predecessor}
\end{align}
The predecessor is either another job or the dummy origin when $j$ is first.

\noindent\textbf{A3. Every job has at most one successor}
\begin{align}
\sum_{h\in\mathcal{J}_{m}\setminus\{j\}}y_{jh}\leq1,
\qquad
\forall m\in\mathcal{M}_{\kappa},\quad
\forall j\in\mathcal{J}_{m}.
\label{eq:A3-successor}
\end{align}

Constraints A1--A3 form one path starting at $0_m$ and covering all jobs on
machine $m$. A disconnected cycle cannot satisfy the timing constraints below
when every active job has $p_j>0$; following the cycle would require a job to
finish before its own start. Thus, no separate subtour-elimination constraints
are required.

\noindent\textbf{A4. A continued job is first}
\begin{align}
y_{0_mj}=1,
\qquad
\forall j\in\mathcal{J}_{m}:\hat c_j=1.
\label{eq:A4-continuation}
\end{align}

\noindent\textbf{A5. The selected end state is last}

Let $succ^{\mathcal{E}}_{m}(j)=\{h\in\mathcal{J}_{m}:(j,h)\in\mathcal{E}\}$ be
the same-machine binding successors of $j$.
\begin{align}
\sum_{h\in\mathcal{J}_{m}\setminus\{j\}}y_{jh}=0,
\qquad
\forall j\in\mathcal{J}_{m}:\hat e_j=1,\ \hat c_j=0,\
succ^{\mathcal{E}}_{m}(j)=\emptyset.
\label{eq:A5-end-state}
\end{align}
A5 is skipped for an end-state job in two situations where it would conflict
with a harder requirement. First, when $\hat c_j=1$ too: the daily model
picked the same job as both continuation and end state, so A4 already forces
it first, and a literal A5 would additionally force it last, leaving every
other job on the machine without a valid position. Second, when $j$ has a
same-machine binding successor $h$ ($succ^{\mathcal{E}}_{m}(j)\neq\emptyset$):
C1 already forces $h$ to run after $j$, so a literal A5 would again forbid
$j$ from having any successor. In both cases A4 and C1 win because they
reflect a real physical setup state or a genuine production dependency, while
$\hat e_j$ is only the daily model's aggregate preference. The job the
sequencing solution actually places last becomes $0_m$ for the next shift,
which may then differ from the daily model's original end-state choice.

\subsubsection{B. Timing}

\noindent\textbf{B1. Completion time}

The selected incoming arc simultaneously identifies the predecessor of $j$
and its setup time:
\begin{align}
C_j
=
B_j+p_j
+
\sum_{i\in\mathcal{N}_{m}\setminus\{j\}}d_{ij}y_{ij},
\qquad
\forall m\in\mathcal{M}_{\kappa},\quad
\forall j\in\mathcal{J}_{m}.
\label{eq:B1-completion}
\end{align}

\noindent\textbf{B2. Machine succession}
\begin{align}
B_j
\geq
C_i-M(1-y_{ij}),
\qquad
\forall m\in\mathcal{M}_{\kappa},\quad
\forall i,j\in\mathcal{J}_{m}:i\neq j.
\label{eq:B2-machine-succession}
\end{align}
If $y_{ij}=1$, the setup of $j$ starts only after job $i$ is complete.

\noindent\textbf{B3. Component makespan}
\begin{align}
C^{\max}_{\kappa}\geq C_j,
\qquad
\forall j\in\mathcal{J}_{\kappa}.
\label{eq:B3-makespan}
\end{align}

\subsubsection{C. Binding precedence}

\noindent\textbf{C1. Within-shift production precedence}
\begin{align}
B_j\geq C_i,
\qquad
\forall(i,j)\in\mathcal{E}.
\label{eq:C1-precedence}
\end{align}
Because $B_j$ is the beginning of setup, this is a non-anticipatory setup
assumption: neither setup nor processing of $j$ may begin before its binding
predecessor is complete.

\subsection{Interpretation and scope}

\begin{itemize}

\item The daily and sequencing models are solved consecutively, never jointly.
The daily solution fixes jobs, quantities, machines, and binding precedences;
the present model only orders the fixed jobs.

\item Shifts are solved chronologically. The last scheduled job on a machine
determines the dummy-origin state $0_m$ used in its next active shift. An idle
shift does not erase the carried state.

\item A job is atomic and non-preemptive. Its setup and processing are
consecutive and occupy the same machine.

\item The shift duration does not bound $C^{\max}_{\kappa}$. Completion beyond
the available minutes is an output of the model, not an infeasibility.

\item For a component containing a single machine and no binding precedence,
total processing time is fixed, so the objective reduces to choosing the
job order that starts at $0_m$ and minimizes total setup time.

\end{itemize}

\end{document}
